\documentclass[c5size,a4paper,nofonts]{ctexart}
\usepackage[utf8]{inputenc}
\def\tjf{{\tt{田劲锋}}}
\def\titlec{页式存储管理的页面置换算法模拟}

\usepackage[a4paper,margin=2.2cm]{geometry} % 页面设置
\usepackage[unicode,breaklinks=true,
colorlinks=true,linkcolor=black,anchorcolor=black,citecolor=black,urlcolor=black,
pdftitle={\titlec},pdfauthor={\tjf}]{hyperref}
\usepackage{latexsym,amsmath,amssymb,bm}
\usepackage{graphicx}
\usepackage{xcolor}
\usepackage{fancyvrb}
\DefineShortVerb{\|}
\fvset{frame=single}

\setmainfont{Times New Roman}
\setCJKmainfont[BoldFont={SimHei}]{SimSun}  % 主要字体：宋体、黑体
\setCJKsansfont[BoldFont={STZhongsong}]{STFangsong} % 次要字体：仿宋、中宋
\setCJKmonofont{KFKai} % 等宽字体：楷体

\usepackage{listings} % 排版代码
\lstset{language=C,
  numbers=left,
  numberstyle=\small,
  basicstyle=\tt\color{black},
  commentstyle=\color{gray},
  keywordstyle=\bfseries\color{blue},
  stringstyle=\color{teal},
  showstringspaces=false,
  frame=tb,
  extendedchars=false,
%     % identifierstyle=\sf,
%     breaklines=true,
  texcl=true,
  mathescape=true,
  escapeinside=` `,
%     % texcl=true,  % 为true时，注释内容被当做LaTeX代码来处理
}


\CJKsetecglue{\hspace{0.1em}}
\renewcommand\CJKglue{\hskip -0.3pt plus 0.08\baselineskip}
\frenchspacing
\widowpenalty=10000
\linespread{1.2} % 行距

\CTEXsetup[name={实验,},number={\arabic{part}}]{part}
\CTEXsetup[name={,.},format={\large}]{section}

\usepackage[inline]{enumitem} % 调整列表样式
\setlist{noitemsep}
\setlist[itemize]{topsep=0pt,partopsep=0pt,itemsep=0pt,parsep=0pt}
\setlist[enumerate]{topsep=0pt,partopsep=0pt,itemsep=0pt,parsep=0pt}
\setlist[enumerate,1]{label={\arabic*.}}
\setlist[enumerate,2]{label={(\arabic*)}}

%\makeindex
\pagestyle{plain}

\newcommand{\image}[3][height=10cm]{%\begin{minipage}[t]{0.5\textwidth}
    \centering
    \includegraphics[#1]{#2}
    \caption{#3}
    \label{fig:#3}
%\end{minipage}
}

\linespread{1.1} % 行距
\fvset{frame=none}
\UndefineShortVerb{\|}

\begin{document}

%%%% 开始 %%%%

\setcounter{part}{4}
\begin{titlepage}

\begin{center}

\includegraphics[height=1cm]{image/haut.png}

\vspace*{1cm}
{\liti\fontsize{48pt}{50pt}{课\quad 程\quad 设\quad 计}}

\vspace*{4cm}
{\fontsize{36}{80}\sf\bfseries \titlec}

\vspace*{1cm}
{\fontsize{30}{70}\sf\bfseries \titlee}

\vfill
{\large
\newcommand{\ctline}[2]{\makebox[6em][s]{\bf #1}：\underline{\makebox[14em][c]{\qquad #2\qquad}}\\}
\ctline{课程设计名称}{数据结构课程设计}
\ctline{专业班级}{计算机 1303 班}
\ctline{学生姓名}{\tjf}
\ctline{学号}{201316920311}
\ctline{指导教师}{白\quad 浩}
\ctline{课程设计时间}{\today}
}

\end{center}

\end{titlepage}


\iffalse
\section{实验目的}
通过请求页式存储管理中页面置换算法模拟设计，了解虚拟存储技术的特点，掌握请求页式存储管理的页面置换算法。

\section{实验环境}
装有操作系统Windows XP和开发工具VC++6.0，内存在256M以上的微机；

或者：装有Linux(Fedora 7)操作系统和gcc编译器，内存在256M以上的微机。

\section{实验内容}
\begin{enumerate}[label={(\arabic*)}]
\item 通过随机数产生一个指令序列，共320条指令。指令的地址按下述原则生成：
\begin{enumerate}[label={\textcircled{\small\arabic*}}]
\item 50\%的指令是顺序执行的；
\item 25\%的指令是均匀分布在前地址部分；
\item 25\%的指令是均匀分布在后地址部分；
\end{enumerate}

具体的实施方法是：
\begin{enumerate}[label={\textcircled{\small\arabic*}}]
\item 在$[0,319]$的指令地址之间随机选取一起点$m$；
\item 顺序执行一条指令，即执行地址为$m+1$的指令；
\item 在前地址$[0,m+1]$中随机选取一条指令并执行，该指令的地址为$m'$；
\item 顺序执行一条指令，其地址为$m'+1$的指令；
\item 在后地址$[m'+2,319]$中随机选取一条指令并执行；
\item 重复上述步骤\textcircled{\small1}～\textcircled{\small5}，直到执行320次指令。
\end{enumerate}

\item 将指令序列变换为页地址流
\begin{enumerate}[label={\textcircled{\small\arabic*}}]
\item 设页面大小为1K；
\item 分配内存容量为4K到32K；
\item 用户虚存容量为32K。
\end{enumerate}

在用户虚存中，按每K存放10条指令排列虚存地址，即320条指令在虚存中的存放方式为：\\
第0条～第9条指令为第0页（对应虚存地址为$[0,9]$）；\\
第10条～第19条指令为第1页（对应虚存地址为$[10,19]$）；\\
……\\
……\\
第310条～第319条指令为第31页（对应虚存地址为$[310,319]$）。\\
按以上方式，用户指令可组成32页。
\item 计算先进先出（FIFO）算法或最近最少使用（LRU）算法在不同内存容量下的命中率。

其中，$\text{命中率}=1-\text{页面失效次数}/\text{页地址流长度}$
\end{enumerate}

\section{实验要求}
\begin{enumerate}[label={(\arabic*)}]
\item 将FIFO或者LRU算法的源程序及程序执行结果写入实验报告；
\item 将FIFO和LRU算法的工作机理写入实验报告。
\end{enumerate}

\fi

\section{实验步骤}

\begin{enumerate}

\item 以下是{\tt page.c}的源代码，注释已详细给出：

{\small\linespread{1}\lstinputlisting[caption={\tt page.c}]{exp5/page.c}}

编译并执行该程序：

\begin{Verbatim}[frame=single]
$ cc -Wall page.c -o page
$ ./page > 1
\end{Verbatim}

得到输出结果如下：

\VerbatimInput[fontsize=\small,frame=lines]{exp5/1}

可以看到两种算法的差距并不是很大，相对的来说，LRU算法在内存较大的情况下会有比FIFO更好一些的效果。

\item {\bf 先进先出}算法（First-In, First-Out, FIFO）是一个经典的队列实现，它也是一种低开销的页面置换算法。我们维护一个链表，链表中记录了所有位于内存中的虚拟页面，从链表的排列顺序看，链表头部页面的驻留时间最长，链表尾部页面的驻留时间最短。当发生一个缺页中断时，将链表头部的页面淘汰出局，并把新的页面添加到链表尾部。FIFO算法的性能不是很好，被它淘汰出去的页面可能是经常要访问的页面。基于这个原因，FIFO算法很少被单独使用。

{\bf 最近最久未使用}算法（Least Recently Used, LRU）的基本思路是，当一个缺页中断发生时，从内存中选择最久未被使用的那个页面，把它淘汰出局。这种算法实质上是对最优页面置换算法的一个近似，理论依据就是程序的局部性原理。LRU算法的策略就是根据程序的局部性原理，利用过去的、已知的页面访问情况，来预测将来的情况。LRU的实现开销比较大，改程序中的模拟实际上效率是比较差的。一种好的解决方案是{\bf 最不经常使用}算法（Not Frequently Used, NFU）的变体——老化算法，具体可以参考《操作系统设计与实现（第三版）》第273--275页的内容。

\end{enumerate}

\end{document}
